\uuid{IsUy}
\titre{Ce que représente une densité}
\auteur{Maxime Nguyen}
\organisation{AMSCC}
\datecreate{2026-11-24}
\niveau{L2}
\module{Probabilité}
\chapitre{Probabilité continue}
\sousChapitre{Densité de probabilité}
\difficulte{1}
\theme{densité, probabilité d'un point}

\contenu{
\begin{question}
Soit $X$ une variable aléatoire de densité $f$. Quelles affirmations sont vraies ?
\begin{answers}
\bad{Pour tout réel $x$, $f(x) = \mathbb{P}(X = x)$.}
\bad{Pour tout réel $x$, $f(x) \leq 1$.}
\good{$\displaystyle\int_{-\infty}^{+\infty} f(x)\,\mathrm{d}x = 1$.}
\good{Pour tout réel $x$, $\mathbb{P}(X = x) = 0$.}
\bad{$f$ est continue sur $\mathbb{R}$.}
\end{answers}
\begin{explanations}
A est faux : $\mathbb{P}(X = x) = \int_x^x f = 0$, alors que $f(x)$ peut être non nul.

B est faux : la fonction qui vaut $2$ sur $[0;1/2]$ et $0$ ailleurs est une densité (positive, d'aire $2 \times \frac12 = 1$), et elle dépasse $1$.

D est vrai pour toute variable à densité.

E est faux : la densité de la loi uniforme sur $[a;b]$, qui vaut $\frac{1}{b-a}$ sur $[a;b]$ et $0$ ailleurs, est discontinue en $a$ et en $b$.

Contre-sens visés : confondre la densité avec une probabilité (une densité
n'est pas une probabilité ; c'est l'aire sous sa courbe qui en est une) ;
croire qu'une variable « continue » a forcément une densité continue.
\end{explanations}
\end{question}
}
